\documentclass[UTF8,aspectratio=169]{ctexbeamer}
\usetheme{Madrid}
\usecolortheme{whale}
\usepackage{amsmath,amssymb}
\usepackage{booktabs}
\usepackage{graphicx}

\title[ECT ICLR 2027 汇报]{ECT 项目 ICLR 2027 战略汇报}
\subtitle{实测诊断:从"机制+方法"到"诚实负结果"}
\author{Role C --- 理论负责人}
\date{2026-08-10}

\begin{document}

\begin{frame}
\titlepage
\end{frame}

\begin{frame}{战略轨迹}
\begin{center}
\begin{tabular}{@{}ll@{}}
\toprule
阶段 & 结论 \\
\midrule
机制+方法($40$--$50\%$) & 被 $g{=}1.3$ 赢 FID 证伪 \\
诚实诊断($20$--$30\%$) & 被"良性"非结果质疑 \\
诊断信号转向($30$--$35\%$) & 被 P1 附庸性检验证伪 \\
\bottomrule
\end{tabular}
\end{center}
\vspace{1em}
\textbf{最终定位:workshop 级诚实负结果,不投 ICLR 主会。}
\end{frame}

\begin{frame}{实测诊断证据(6 个实验)}
\begin{center}
\small
\begin{tabular}{@{}lll@{}}
\toprule
环节 & 实验 & 结论 \\
\midrule
存在 & moment-memory & 标量 null 被证伪($h$ $1.001$ vs.\ $0.837$) \\
有结构 & E5 & 深度趋势 + 剂量响应,非噪声 \\
无害 & E6 + 干净 FID & 与 FID 反相关 $r{=}{-}0.99$ \\
状态条件 & E7 & $R_{\mathrm{grad}}$ 随 $n_K$ 增长 \\
稳健 & E11 & 跨 3 seed + 支撑阈值 \\
普遍 & E3 & 梯度残差跨 optimizer 复现 \\
\bottomrule
\end{tabular}
\end{center}
\vspace{0.5em}
\small 关键数字:$a_K^*{=}0.761$、$R_{\mathrm{grad}}{=}5.85\%$、$R_{\mathrm{opt}}{=}8.57\%$、$H_K{=}R_{\mathrm{opt}}$ 精确。
\end{frame}

\begin{frame}{P1 决定性结果:残差附庸于标量 gap}
\begin{center}
\begin{tabular}{@{}ll@{}}
\toprule
相关 & 值 \\
\midrule
$\mathrm{corr}(\text{残差},\text{FID})$ & $-0.817$ \\
$\mathrm{corr}(|g-1|,\text{FID})$ & $-0.899$ \\
$\mathrm{corr}(\text{残差},|g-1|)$ & $+0.938$ \\
\textbf{偏相关(残差, FID $|$ $|g-1|$)} & \textbf{$+0.168\approx 0$} \\
\bottomrule
\end{tabular}
\end{center}
\vspace{0.5em}
残差几乎由 $|g-1|$ 决定,在标量 gap 之外\textbf{不增加对 FID 的预测力}。
诊断信号论题不成立。
\end{frame}

\begin{frame}{诊断规划(来自 ICLR2027\_STRATEGY)}
\textbf{诊断论文主张:} 非标量梯度内容在少步训练中系统性存在、有结构、状态条件、
普遍、\emph{良性}——标量 moment-memory 恒等式($h{=}1.001$)被实测($h{=}0.837$)
证伪。此点未被 2025--26 标量 scale-invariance 文献覆盖。

\textbf{证据链(已实测):} 存在(moment-memory)→ 有结构(E5)→ 无害(E6 + 干净 FID)
→ 状态条件(E7)→ 稳健(E11)→ 普遍(E3)。

\textbf{边界(四重实测证据):} P1 附庸性(偏相关 $\approx 0$)+ Arm C scale-matched
control($g{=}1.3$ 改善被消除 84--100\%)+ 梯度级未来预测阴性(早期 $R_{\mathrm{grad}}$)+
优化器级未来预测阴性(早期 $R_{\mathrm{opt}}$ 不预测未来 FID)。残差是真实但无后果的
现象。\textbf{不声称}诊断信号、\textbf{不声称}有害/可修正、\textbf{不投} ICLR 主会。
\textbf{Go/No-Go 定论:No-Go for ICLR 主会机制论题。} 诚实的 workshop 级结构刻画是最终可发表资产。

\textbf{剩余规划:} 更新定位为 workshop 级诚实负结果;投 ``I Can't Believe It's
Not Better'' workshop;可选补强第二数据集挣"普遍";不可行项降级 future work。
\end{frame}

\begin{frame}{最终定位}
\begin{block}{诚实结论}
非标量梯度残差存在、有结构、状态条件、普遍、\emph{良性}——但\emph{附庸于标量
gap}。它是少步训练的真实特征,但不是质量诊断信号,也不有害。
\end{block}
\begin{alertblock}{论文定位}
\textbf{workshop 级诚实负结果,不投 ICLR 主会。}
\end{alertblock}
\end{frame}

\begin{frame}{交付物与待办}
\textbf{交付物:}
\begin{itemize}
  \item 诊断论文:`docs/ICLR2027\_DIAGNOSIS\_PAPER.{md,tex,pdf}`
  \item 战略文档:`docs/ICLR2027\_STRATEGY.{tex,pdf,md}`
  \item 对抗性审查:`docs/ICLR2027\_PLAN\_REVIEW.md`
  \item 实验结果(7 份):`analysis/g13\_vs\_g10\_fid\_result.md` 等
  \item Git:分支 `role-c/g13-vs-g10-fid-0809`,23 commit 领先 main,未合并
\end{itemize}
\textbf{待办:}
\begin{enumerate}
  \item 本地 push 分支(勿推 main)
  \item 更新论文定位为 workshop 级诚实负结果
  \item 可选:投 ``I Can't Believe It's Not Better'' workshop
\end{enumerate}
\end{frame}

\begin{frame}{谢谢}
\begin{center}
\Large 谢谢聆听
\end{center}
\end{frame}

\end{document}
